\documentclass{llncs}

\usepackage{graphicx}
\usepackage{grffile}
\usepackage{amsmath}
\usepackage{enumitem}

\begin{document}

\title{Motion Stopping and Stable Deceleration of Legged Robots}

\author{Krishna Sikheriya (IIT2023139) \and Grish Gautam (IIT2023131) \and
Nitya Bhavsar (IIT2023140) \and Priyam Jyoti Chakrabarty (IIT2023147) \and
Tavish Chawla (IIT2023150) \and Biswajit Dash (MRM2025003)}

\institute{Group 2, Foundations of Robotics (Mathematical Foundation of Robotics), BTech IT}

\maketitle

\begin{abstract}
A legged robot that is walking or running carries forward momentum, and simply locking its legs
will not bring it to a safe stop: the body keeps moving and the robot falls over. This report
covers the first phase of our project, where we study the Capture Point method, which gives the
exact spot on the ground where the next foot must land to cancel out this momentum and bring the
robot to rest. We implement the method on a simplified point-mass model of a legged robot,
compare it against a naive controller that simply locks the legs, and measure the fastest speed
from which the robot can still stop safely. The report also summarizes the literature behind this
approach and the work completed so far.
\keywords{legged robots, capture point, inverted pendulum, balance control, stopping stability}
\end{abstract}

\section{Team Details}
\label{sec:team}

\textbf{Assigned Group Number:} 2

\begin{itemize}[leftmargin=1.5cm]
    \item Krishna Sikheriya (IIT2023139)
    \item Grish Gautam (IIT2023131)
    \item Nitya Bhavsar (IIT2023140)
    \item Priyam Jyoti Chakrabarty (IIT2023147)
    \item Tavish Chawla (IIT2023150)
    \item Biswajit Dash (MRM2025003)
\end{itemize}

\section{Problem Statement}
\label{sec:problem}

When a legged robot is walking or running and needs to stop, it cannot simply freeze its legs.
The body still has forward momentum, and if the legs lock in place, that momentum carries the
body past its feet and the robot falls over, much like a person tripping while running. To stop
safely, the robot has to take one or more quick steps and place its foot at the correct spot on
the ground so that the momentum is cancelled out and the body comes to rest without falling.

Our project looks at the math and control logic behind finding that correct foot placement. We
are working with a simplified model of a legged robot, where the body is treated as a point mass
balanced on a single leg. Using this model, we compute a value called the Capture Point, which
tells us exactly where the next foot needs to land to bring the robot to a stop. We compare this
against a simple, naive controller that just holds the leg in place once told to stop, to clearly
show why that approach fails.

We are also studying how well this stopping method holds up under different conditions, such as
different starting speeds, pushes from the side, and different ground surfaces like slopes or
slippery patches. Since our team does not have access to a physical robot or prior robotics
background, all of this work is done through software simulation, first as a simplified math
model and later as a 3D simulation. The end goal is to understand and demonstrate a working
method for stable stopping that could, in principle, be used on a real legged robot.

\section{Literature Review}
\label{sec:litreview}

Our approach draws on ten papers, two of them older but still central to almost everything
published on this topic since.

The core of the project comes from Koolen et al. (2012), who worked out the idea of the Capture
Point and N-step capturability, a formal way of asking whether a legged robot can come to rest
without falling within N steps. Chen et al. (2023) later carried this over to four-legged robots
by pairing the Linear Inverted Pendulum model with a model predictive control scheme, and it is
essentially this LIP model that Phase 1 is built around.

A few of the more recent papers are less about the mechanics of stopping and more about deciding
when to stop. He et al. (2024), which won the Best Paper award at RSS that year, trained a
quadruped with two separate policies, one for fast running and one purely for recovery, and
switched between them using a value function derived from reachability analysis. Sun et al.
(2024) took a simpler route to a similar end, using convex optimization to continuously check
whether a fall was becoming likely and triggering a dedicated stopping controller before it was
too late. Both reinforced the same point for us: stopping works better as its own explicit mode
with a clear switching condition than as something added onto the normal walking controller
afterwards, which is roughly how our controller is structured, using a plain stop flag rather
than anything learned.

Two other papers shaped how we test and measure the controller rather than how we build it. van
Marum et al. (2024) proposed a walk/stand command interface along with concrete benchmarking
metrics, such as how far a robot drifts before it stops and how much energy it uses, and we
borrowed these loosely as a way of judging our own results. Weng et al. (2024) described a
standardized mechanical device for pushing a robot repeatably during testing; we cannot build a
physical version of it, but the same idea carries over directly into a simulated, precisely
timed push once we get to Phase 2.

Subburaman et al. (2023) surveyed the wider field of humanoid falling and organized it into three
stages: predicting a fall, controlling the stop or fall itself, and recovering afterward, and we
have loosely followed that structure in how this project is organized. Thuruthel et al. (2021),
whose title matches our own project almost exactly, argued that learning to stop reliably is a
foundational skill rather than an afterthought, which is the framing we use to motivate the work.

Kim et al. (2024) showed that legged robot controllers trained with explicit safety constraints
tend to behave more reliably than ones relying only on shaped rewards, which lines up with our
decision to enforce the Capture Point condition and the leg-reach limit as hard rules instead of
something learned. Han et al. (2022) used a more detailed 3D Spring-Loaded Inverted Pendulum
model, with an added energy-compensation term, to handle sideways impacts during running; this is
a reasonable next step if time allows once the core controller is fully working.

Of the ten, Koolen et al. (2012) and Chen et al. (2023) matter most directly, since the Linear
Inverted Pendulum model and Capture Point formula used throughout this project come from them.
The rest inform how we test, evaluate, and frame the work rather than the underlying math.

\section{Progress and Implementation}
\label{sec:progress}

\subsection{Repository Setup}
\label{sec:repo}

We created a GitHub repository for the project and set up a Python environment using a virtual
environment (\texttt{venv}) with \texttt{numpy}, \texttt{scipy}, \texttt{matplotlib}, and
\texttt{pandas} for Phase 1. \texttt{pybullet} and \texttt{imageio} are listed for Phase 2 and
will be installed later. A \texttt{requirements.txt} file is included so any team member can set
up the same environment with one command. The repository is organized as follows:

\begin{verbatim}
Phase 1/
    lip_model.py          - LIP dynamics
    capture_point.py       - Capture point formula and controller
    naive_controller.py    - Baseline "locked leg" controller
    run_experiments.py     - Runs the speed sweep and generates all plots
    plots/                 - Generated graphs
    results.csv            - Generated results table
docs/
    literature_review.md   - Full write-up of the papers above
README.md
requirements.txt
\end{verbatim}

\subsection{Core Math Implemented}
\label{sec:math}

We modeled the robot as a point mass at a fixed height above the ground, balanced over one foot.
The horizontal acceleration of this point mass follows:
\begin{equation}
\ddot{x} = \frac{g}{z_0}(x - p)
\label{eq:lip}
\end{equation}
where $x$ is the position of the body, $p$ is the position of the foot, $z_0$ is the fixed body
height, and $g$ is gravity. This equation is unstable on its own: if the body drifts even
slightly away from being directly above the foot, it keeps drifting further away. This instability
is exactly why the foot placement matters.

To stop the robot safely, we compute the Capture Point:
\begin{equation}
T_c = \sqrt{z_0/g}, \qquad x_{capture} = x + \dot{x}\,T_c
\label{eq:capture}
\end{equation}
This tells us where the next foot must land so that the body's velocity settles to zero instead of
growing. We also extended this same formula so it can describe steady walking at a chosen speed,
not just stopping, by using the same idea with a target speed instead of zero.

\subsection{Code and Controllers}
\label{sec:code}

We implemented three main pieces of code inside \texttt{Phase 1/}:

\begin{itemize}[leftmargin=1.5cm]
    \item \textbf{lip\_model.py} -- simulates the pendulum equation above using simple Euler
    integration (small time steps).
    \item \textbf{capture\_point.py} -- computes the Capture Point at each step and decides where
    the foot should go, limited by how far a leg can physically reach in one step.
    \item \textbf{naive\_controller.py} -- a baseline controller that just keeps the foot fixed
    once told to stop, so we can clearly show what happens without the Capture Point method.
\end{itemize}

Both controllers walk in exactly the same way before the stop command is given, so the only
difference being tested is what happens at the moment of stopping.

\subsection{Experiments and Results}
\label{sec:results}

Using \texttt{run\_experiments.py}, we tested both controllers across starting speeds from 0.2
m/s to 4.0 m/s. The results were exactly what we expected:

\begin{itemize}[leftmargin=1.5cm]
    \item The naive controller falls at every single speed we tested, since holding the leg still
    can never fix the instability once the body starts moving.
    \item The Capture Point controller stops the robot cleanly at every speed up to about 3.0
    m/s, by taking one or more corrective steps as needed.
    \item If we only allow the robot a single corrective step, the highest speed it can stop from
    is about 1.93 m/s (this matches the formula max leg reach divided by $T_c$). Allowing a few
    extra steps pushes this limit up to around 3.0 m/s, after which the controller runs out of
    steps and eventually falls too.
\end{itemize}

For each of these tests, the script generates position-vs-time graphs, velocity-vs-time graphs, a
graph comparing the Capture Point to where the foot actually lands, and a summary graph showing
the stopping outcome (stopped, ran out of steps, or fell) across all tested speeds; these are
saved automatically under \texttt{Phase 1/plots/}. Figures~\ref{fig:stability}--\ref{fig:capturevfoot}
show the main results.

\begin{figure}[h]
    \centering
    \includegraphics[width=0.85\textwidth]{Phase 1/plots/stability_limit.png}
    \caption{Stopping outcome vs. starting speed, for both controllers. The naive controller
    (orange) falls at every speed tested. The Capture Point controller (blue) stops successfully
    up to about 3.0 m/s, briefly fails to converge just past that point, and then falls as well,
    giving a clear empirical stability limit.}
    \label{fig:stability}
\end{figure}

\begin{figure}[h]
    \centering
    \includegraphics[width=0.85\textwidth]{Phase 1/plots/trajectory_v1.8.png}
    \caption{Position and velocity over time at a starting speed of 1.8 m/s, comparing both
    controllers. Right after the stop command (dotted line), the naive controller's velocity
    grows instead of decaying, the instability of Equation~\eqref{eq:lip} taking over, while the
    Capture Point controller's velocity decays smoothly to zero.}
    \label{fig:trajectory}
\end{figure}

\begin{figure}[h]
    \centering
    \includegraphics[width=0.85\textwidth]{Phase 1/plots/capture_vs_foot_v2.8.png}
    \caption{Capture point vs. actual foot placement at a starting speed of 2.8 m/s, for the
    Capture Point controller. The foot placement catches up to the capture point over several
    corrective steps rather than in one, illustrating the N-step capturability described by
    Koolen et al. (2012) and explaining why the empirical stopping limit (about 3.0 m/s) exceeds
    the theoretical single-step limit (about 1.93 m/s).}
    \label{fig:capturevfoot}
\end{figure}

\subsection{Future Work}
\label{sec:future}

This covers the pre-midsem part of the project, which is the math model and the Capture Point
controller. The next phase, planned for after the midsem, is to apply this same controller to a
3D robot model in PyBullet, test it against pushes and different ground surfaces, and record a
short demo video comparing the naive and Capture Point controllers.

\clearpage
\begin{thebibliography}{10}

\bibitem{koolen2012}
Koolen, T., de Boer, T., Rebula, J., Goswami, A., Pratt, J.: Capturability-based analysis and
control of legged locomotion, Part 1: Theory and application to three simple gait models.
International Journal of Robotics Research (2012). https://doi.org/10.1177/0278364912452673

\bibitem{chen2023}
Chen, H., Hong, Z., Yang, S., Wensing, P.M., Zhang, W.: Quadruped capturability and push
recovery via a switched-systems characterization of dynamic balance. IEEE Transactions on
Robotics (2023). https://arxiv.org/abs/2201.11928

\bibitem{he2024}
He, T., Zhang, C., Xiao, W., He, G., Liu, C., Shi, G.: Agile but safe: Learning collision-free
high-speed legged locomotion. Robotics: Science and Systems (2024).
https://arxiv.org/abs/2401.17583

\bibitem{sun2024}
Sun, H., Yang, J., Jia, Y., Zhang, C., Yu, X., Wang, C.: Fall prediction, control, and recovery
of quadruped robots. ISA Transactions (2024). https://doi.org/10.1016/j.isatra.2024.05.039

\bibitem{vanmarum2024}
van Marum, B., Shrestha, A., Duan, H., Dugar, P., Dao, J., Fern, A.: Revisiting reward design
and evaluation for robust humanoid standing and walking. IEEE/RSJ International Conference on
Intelligent Robots and Systems (2024). https://arxiv.org/abs/2404.19173

\bibitem{weng2024}
Weng, B., Castillo, G.A., Kang, Y.S., Hereid, A.: Towards standardized disturbance rejection
testing of legged robot locomotion with linear impactor. IEEE International Conference on
Robotics and Automation (2024). https://arxiv.org/abs/2308.14636

\bibitem{subburaman2023}
Subburaman, R., Kanoulas, D., Tsagarakis, N., Lee, J.: A survey on control of humanoid fall
over. Robotics and Autonomous Systems (2023). https://doi.org/10.1016/j.robot.2023.104443

\bibitem{thuruthel2021}
Thuruthel, T.G., Picardi, G., Iida, F., Laschi, C., Calisti, M.: Learning to stop: A unifying
principle for legged locomotion in varying environments. Royal Society Open Science (2021).
https://doi.org/10.1098/rsos.210223

\bibitem{kim2024}
Kim, Y., Oh, H., Lee, J., Choi, J., Ji, G., Jung, M., Youm, D., Hwangbo, J.: Not only rewards
but also constraints: Applications on legged robot locomotion. IEEE Transactions on Robotics
(2024). https://arxiv.org/abs/2308.12517

\bibitem{han2022}
Han, B., Yi, H., Xu, Z., Yang, X., Luo, X.: 3D-SLIP model based dynamic stability strategy for
legged robots with impact disturbance rejection. Scientific Reports (2022).
https://doi.org/10.1038/s41598-022-09937-9

\end{thebibliography}

\end{document}
